\solchapter{continuum}

\sol{continuum}{10} Write the total energy as a function of crack half-length, $U(a) = 4a\gamma - \pi a^2\sigma^2/E$. Where is its maximum, and what happens to a crack longer than that?

\sol{continuum}{11} The modes with frequency below $F$ correspond to lattice points $(m,n)$ inside a circle. For large $F$, count lattice points by area.

\sol{continuum}{12} Divide each measured frequency by that of the $(1,1)$ mode and compare with $1:1.5:2:2.5$. For the $(0,n)$ modes, think about what they do to the volume of air in the kettle, and where on the head the drum is struck.

\solchapter{complex}

\sol{complex}{9} Solve the ODE to get $k_i(t)$ in terms of the arrival time $t_i$. A node has degree greater than $k$ if it arrived early enough; arrival times are uniform.

\sol{complex}{10} Pairwise exchanges conserving a total are like collisions conserving energy. What distribution maximizes the entropy of a fixed total shared among many agents?

\sol{complex}{11} For the code, keep a stack of sites with four or more grains and topple them until none remain. For the second part, ask what happens to the average slope when avalanches are too rare, and when they are too large.

\solchapter{plasma}

\sol{plasma}{9} Compute $\omega/k$ and $\dd\omega/\dd k$ directly. Then set $E = \hbar\omega$, $p = \hbar k$ and compare with $E^2 = (mc^2)^2 + (pc)^2$.

\sol{plasma}{10} The grad-$B$ drift has opposite signs for opposite charges. What field does the resulting charge separation create, and what drift does that field produce, for every particle?

\sol{plasma}{11} Go to the wave's frame. Particles moving slightly slower or faster than the wave are reflected by its potential hills. Which group gains energy, and which loses it?

\solchapter{amo}

\sol{amo}{11} Find the steady state of $\dot N_2 = BI(N_1-N_2) - AN_2$. When $N_1 = N_2$, what do absorption and stimulated emission do to each other?

\sol{amo}{12} Average $V(\vb r+\delta\vb r)$ over random isotropic jitter. The first-order term averages to zero; the second-order term involves $\nabla^2V$, which for the Coulomb potential is a delta function.

\solchapter{condmat}

\sol{condmat}{11} Write the two coupled equations for the amplitudes $a$ and $b$ on the two sublattices and multiply them together. Where in $k$ is $\abs{t_1+t_2e^{ikd}}$ smallest?

\sol{condmat}{12} Faraday's law turns a change of flux into an EMF around the cylinder; the Hall conductance turns that EMF into a current along it. Integrate over the time it takes to add one flux quantum.
